\chapter{Sprint 63 --- Cụm F khối A: tách 7 phân hệ khỏi portal + controller mỏng, F6 $\to$ $\star$FR-A (2026-09-25 $\to$ 09-26)}

\emph{Chapter này ghi lại nhánh sau cổng của cụm F. Cổng F0 $\to$ $\star$FR-A3 (chapter trước) đã dọn CSS,
menu và test của khung. Khối A làm nốt phần còn lại của ADR-027 (chapter 74): đưa nghiệp vụ ra khỏi
7 module \texttt{wujia\_portal\_*} và làm mỏng controller. Khối gồm 11 phiên trong hai ngày; mỗi phiên tách
xong đều lên UAT ngay trong ngày. Không có tính năng mới. Người dùng portal và backend không thấy khác gì,
trừ vài chỗ sửa lỗi có từ trước (liệt kê ở cuối). Nhật ký ở \texttt{docs/f-progress.md}; nghiệm thu từng phiên ở
\texttt{docs/f6-} \ldots\ \texttt{f13-acceptance-matrix.md}; hai phiên review ở \texttt{docs/f-review-FR-P.md} và
\texttt{docs/f-review-FR-A.md}.}

\section{Vì sao có khối A}

Review 17/09 đo được: 7 module portal (\texttt{order\_window}, \texttt{info\_request}, \texttt{knowledge},
\texttt{support}, \texttt{notification}, \texttt{exam}, \texttt{return}) chứa cả model, nhóm quyền, sequence,
menu và backend HQ. Hệ quả có hai tầng:
\begin{itemize}[nosep]
\item \textbf{Mobile bị kéo lên portal.} Module mobile cho một phân hệ phải depend module portal chứa model đó.
  Như vậy là phá luật ``hai kênh không depend chéo'' của ADR-027, và không tách repo được.
\item \textbf{Luật nghiệp vụ nằm trong controller.} Cùng một luật (ai thấy bài nào, phiếu nào hợp lệ, tối đa bao
  nhiêu người) được viết lại ở controller, ở Home và ở constraint, và các bản viết đã lệch nhau.
\end{itemize}

Cách làm đã chốt: tách \textbf{từng phân hệ một}, mỗi phiên một lần deploy, và giữ nguyên \texttt{\_name}
(bảng DB không đổi, không migration dữ liệu). Pilot rẻ nhất đi trước (\texttt{order\_window}). Xong pilot có một
phiên review ($\star$FR-P) để quyết có nhân quy trình ra 6 phân hệ còn lại hay không. Xong cả khối có
$\star$FR-A.

\begin{center}
\small
\begin{tabular}{llll}
\toprule
Phiên & Nội dung & Commit & Số đo chốt \\
\midrule
F6 & Controller Đặt hàng mỏng & \texttt{df385d1} & suite 822/0 · mutation 10/10 \\
F7 & Pilot tách \texttt{order\_window} & \texttt{6b2938d} & 42 đổi chủ · HTML 5/5 byte \\
$\star$FR-P & Review pilot, gỡ vỏ trên UAT & \texttt{9d2a606} & UAT 0 lệch · suite 834/0 \\
F8 & \texttt{info\_request} & \texttt{a5cb738} & 80 đổi chủ · 0 lệch thật \\
F9 & \texttt{knowledge} & \texttt{2ba8311} & 122 đổi chủ · số kế KNW- giữ \\
F10 & \texttt{support} (+ fix đính kèm \texttt{daf19f1}) & \texttt{033794c} & 130 đổi chủ · suite 872/0 \\
F11 & \texttt{notification} & \texttt{bfa2bae} & 173 đổi chủ · HTML 117/120 \\
F12 & \texttt{exam} & \texttt{2acbd9f} & 245 đổi chủ · \texttt{check\_layers} 0 \\
F13 & \texttt{return} & \texttt{a138d21} & 287 đổi chủ · suite 912/0 \\
$\star$FR-A & Review toàn khối A & \texttt{2835f8e} & query $\Delta$0 · suite 914/0 \\
\bottomrule
\end{tabular}
\end{center}

\section{F6 --- Controller Đặt hàng mỏng}

\begin{itemize}[nosep]
\item \textbf{Luật số lượng một nguồn}: \texttt{product.\_portal\_qty\_error} (\texttt{wujia\_sale}) được dùng chung
  cho thêm giỏ, sửa số lượng, dòng giỏ và constraint dòng đơn portal. Mã lỗi và câu chữ giữ nguyên.
\item \textbf{Giỏ và gửi đơn về model} \texttt{wujia.portal.cart}. \texttt{action\_submit\_order()} làm trong
  \textbf{một savepoint} các bước: khoá NOWAIT, kiểm khung giờ, kiểm dòng lỗi, huỷ báo giá cũ, tạo đơn, xoá giỏ.
  Controller chỉ còn đổi mã lỗi ra redirect (1047 $\to$ 862 dòng).
\item \textbf{Vá lỗi ẩn}: code cũ để lại đơn nháp mồ côi khi tạo đơn lỗi giữa chừng. Test đặc tả chứng minh lỗi
  trên HEAD trước khi sửa.
\item Nợ FR-A3 trả xong: \texttt{portal\_base} cài một mình giờ tự test được (39 failed + 82 error $\to$ 0/0), nhờ
  \texttt{tests/common.py} cho test quét nhiều module tự bỏ qua khi module chủ chưa cài.
\end{itemize}

\section{F7 và $\star$FR-P --- pilot \texttt{order\_window}}

F7 tạo \texttt{wujia\_order\_window} (L2, depend \texttt{wujia\_sale}). Module này nhận nguyên model khung giờ,
cấu hình fallback và luật chặn đơn portal ngoài giờ, với lớp lỗi riêng \texttt{OrderWindowClosed}. Nhờ có lớp
lỗi riêng, giỏ hàng không còn bắt lỗi bằng cách dò chữ ``khung giờ''. \texttt{pre\_init\_hook} đổi chủ cả
\texttt{ir\_model\_constraint}/\texttt{relation}, chỗ mà khuôn tách Khảo sát trước đó bỏ sót. Công cụ
\texttt{scripts/qa/split\_snapshot.py} (chụp trước/sau, diff có \texttt{--rename}) ra đời ở đây và dùng cho mọi phiên sau.

$\star$FR-P kết luận \textbf{đạt, nhân được}. Vỏ \texttt{wujia\_portal\_order\_window} được gỡ trên UAT qua RPC (lời gọi
ghi duy nhất, có duyệt), đo lại 0 lệch, rồi xoá khỏi repo. Hàm đổi chủ dời về L1
\texttt{wujia\_core/tools/module\_split.py} (\texttt{imd\_names} suy xmlid theo model + \texttt{migrate\_ownership}).
Phiên này cũng bắt được \texttt{deploy.yml} còn \texttt{-i/-u} vỏ. Phần mobile của anh Thái được chốt là mục bàn giao,
không chặn F8.

\section{F8--F13 --- sáu phân hệ theo cùng một khuôn}

Mỗi phiên theo đúng quy trình ở chapter 74 (\S\ref{sec:adr27-xmlid}): đo UAT chỉ-đọc trước, \texttt{git mv} model +
quyền + data sang \texttt{wujia\_<x>}, hook đổi chủ, controller mỏng, snapshot trước/sau trên DB copy, HTML so từng byte,
mutation, rồi deploy một lệnh \texttt{-i wujia\_<x> -u wujia\_portal\_<x>,\ldots}. Module portal còn lại đúng phần kênh:
controller, QWeb, mục điều hướng.

\begin{center}
\footnotesize
\begin{tabularx}{\textwidth}{>{\raggedright\arraybackslash}p{2.5cm}>{\raggedright\arraybackslash}p{4.4cm}>{\raggedright\arraybackslash}X}
\toprule
Module L2 mới & Nhận về & Luật dời từ controller về model \\
\midrule
\texttt{wujia\_info\_request} (F8) & 1 model, rule, sequence \texttt{INF/}; 72 xmlid + 7 ràng buộc &
  \texttt{\_portal\_can\_request}, \texttt{\_portal\_scope\_domain}, \texttt{create\_from\_portal} (savepoint) \\
\texttt{wujia\_knowledge} (F9) & 3 model, cron, \texttt{KNW-}; 107 xmlid + 13 ràng buộc &
  \texttt{\_portal\_visible\_domain} (Home dùng chung), \texttt{\_portal\_search\_domain} \\
\texttt{wujia\_support} (F10) & 2 model, 7 danh mục, \texttt{WJ-TK/}; 113 xmlid + 16 ràng buộc &
  \texttt{create\_from\_portal}, \texttt{\_portal\_reply}, \texttt{\_portal\_get\_attachment} \\
\texttt{wujia\_notification} (F11) & 3 model, 2 nhóm quyền, \texttt{ANN/}; 140 xmlid + 28 ràng buộc &
  \texttt{\_portal\_effective\_domain}, \texttt{\_portal\_unread\_count}, \texttt{\_mark\_read} \\
\texttt{wujia\_exam} (F12) & 5 model, 2 nhóm, 3 sequence; 214 xmlid + 30 ràng buộc &
  \texttt{register\_from\_portal}, \texttt{\_effective\_max\_per\_registration} \\
\texttt{wujia\_return} (F13) & 6 model (kế thừa SO/picking/product), wizard bù, 2 nhóm, 2 sequence; 230 xmlid + 54 ràng buộc &
  \texttt{create\_from\_portal}, \texttt{\_portal\_status\_key}, \texttt{\_portal\_compensation\_view} \\
\bottomrule
\end{tabularx}
\end{center}

Controller các màn giảm theo: yêu cầu thông tin 264 $\to$ 240 dòng · kiến thức 181 $\to$ 154 · hỗ trợ 201 $\to$ 166 ·
thông báo 434 $\to$ 337 · thi 750 $\to$ 579 · đổi trả 609 $\to$ 378. Số route không đổi (103 route / 106 path).
Mọi số kế sequence trên UAT được giữ (INF 1 · KNW 28 · WJ-TK 17 · ANN 20 · WJ-CRS 5 · WJ-EXR 17 · WJ-EXS 6 · RTN 5 · CA 1).

Các bẫy lộ ra trong lúc làm, đều đã ghi vào quy trình ở chapter 74:
\begin{description}[nosep]
\item[Data \texttt{noupdate} (F9).] \texttt{-i} module mới nạp lại bản ghi noupdate lên xmlid vừa đổi chủ. Dòng
  \texttt{number\_next} trong XML reset sequence KNW- từ 438 về 1 trên DB copy, trong khi snapshot cũ vẫn báo 0 lệch.
  Cách sửa: bỏ \texttt{number\_next} khỏi XML, và snapshot đọc thêm \texttt{last\_value} của sequence Postgres.
  F11 là lần đầu dữ liệu UAT lệch XML thật (4 màu loại thông báo); chủ dự án chốt để về theo code.
\item[Mobile tham chiếu xmlid portal (F8, F12).] \texttt{-u} gãy ở module mobile. Theo chốt, đã sửa 5 và 18 dòng
  \texttt{ref} trong module anh Thái. Đo thực tế: Odoo 19 commit sau từng module, nên khi lỗi thì chỉ module mobile dừng lại.
\item[\texttt{groups=} trong arch view (F12).] Tên nhóm thiếu tiền tố module bị bỏ qua mà không báo lỗi, nên 10 nút
  backend bị ẩn với cả Administrator. Snapshot bắt được qua md5 arch.
\item[Test rollback (F8).] \texttt{assertRaises} tự bọc savepoint, nên test vẫn Pass khi code quên savepoint.
  Phải viết bằng \texttt{try/except}.
\end{description}

\section{$\star$FR-A --- review toàn khối A: đạt, khối A khép}

Đo HEAD \texttt{164d8ec} cạnh đối chứng \texttt{9d2a606} (trước F8), hai DB cùng một lần seed:
\begin{itemize}[nosep]
\item \textbf{Tầng}: \texttt{check\_layers} 0 vi phạm phía Dev (R7 còn 2 của anh Thái). 7 module L2 có 0 controller,
  0 static, 0 \texttt{import http}. L3b chỉ giữ 2 ngoại lệ đã chốt: giỏ \texttt{wujia.portal.cart} và AbstractModel
  \texttt{wujia.portal.debt}. \texttt{portal\_base} gọi L2 ở 7 chỗ, đều qua \texttt{hasattr}, không thêm depend.
\item \textbf{Hiệu năng}: số query 17 route $\times$ 4 user giống hệt đối chứng ($\Delta$0). Index đủ trên 6 bảng chính.
  Không có \texttt{search()} trong vòng lặp ở method \texttt{\_portal\_*} mới.
\item \textbf{Hành vi}: HTML 63/68 giống từng byte; 5 chỗ khác đều có chủ đích (nhãn Home F13 + fix công nợ).
  DB trắng chỉ cài 7 L2 chạy 123/0/0. Suite \textbf{914/0/0}. Mutation 7/8.
\item \textbf{Sửa nhỏ theo luật $\star$} (11 file, +16 $-$112, không đổi hành vi): 4 hàm chết, \texttt{is\_read\_by},
  \texttt{ormcache}/import thừa, bẫy \texttt{number\_next} còn sót ở \texttt{INF/}, 14 comment dẫn mã phiên.
\item \textbf{Retest fail \texttt{UI-DATALIST-001}}: \texttt{/portal/debt/payment-history} mobile trước đây không có
  pager và card cao 124px. Sửa: dùng chung lát phân trang với PC và bố cục hàng chuẩn ListCard, còn 10 thẻ cao 104px.
  Đã deploy UAT; sheet chuyển Ready for Retest.
\end{itemize}

\textbf{Smoke chốt sổ (26/09):} lấy DB ở trạng thái FR-P, chạy một lượt đúng lệnh deploy UAT (\texttt{-i/-u} 24 module
+ \texttt{wujia\_portal\_debt}): RC 0, 0 ERROR; cả 7 module L2 cài xong, không module nào kẹt \texttt{to upgrade}.

\section{Gì đã làm (nghiệp vụ)}

Bảy phân hệ (khung giờ đặt hàng, yêu cầu cập nhật thông tin, kiến thức, hỗ trợ, thông báo, thi, đổi trả) giờ là module
nghiệp vụ độc lập. Module nghiệp vụ không biết portal tồn tại, nên backend HQ và sau này app mobile dùng thẳng được mà
không phải cài portal. Luật ``ai thấy gì, gửi được gì'' của từng phân hệ chỉ còn một chỗ, dùng chung cho màn portal,
Home và backend. Vì vậy các lệch giữa Home và màn chi tiết đã được vá (bài hẹn giờ hiện sớm ở Home, số thông báo chưa
đọc ở Home khác chuông). Dữ liệu thật trên UAT (phiếu, số chứng từ, người trong nhóm quyền, menu) không mất dòng nào.

\section{Lý do / trade-off}

\begin{itemize}[nosep]
\item \textbf{Tách một phần, giữ \texttt{\_name}}: không đổi tên bảng nên không có migration dữ liệu. Đổi lại, trong L2 còn
  dư âm tên portal (tham số \texttt{wujia\_portal.*}, field \texttt{portal\_order\_time\_*}, app ``Wujia Portal''). Đổi
  những tên này cần migration riêng, nên để lại.
\item \textbf{Hai model ở lại L3b}: giỏ hàng chỉ tồn tại vì kênh portal; \texttt{wujia.portal.debt} là seam đọc kế toán,
  không có bảng. Cả hai được ghi là ngoại lệ trong ADR-027.
\item \textbf{Home không tách}: \texttt{portal\_base} đọc luật của L2 qua \texttt{hasattr}. Nhờ vậy nó vẫn cài được khi
  thiếu một phân hệ (có test giả lập module tắt), và luật cấm thêm depend vẫn giữ.
\item \textbf{Mỗi phân hệ một lần deploy}: bảy lần đổi chủ xmlid trên dữ liệu thật mà gộp một lệnh thì hỏng sẽ không
  bisect được. Tách lẻ từng phiên tốn thêm thời gian đo nhưng mỗi lần deploy đều 0 lệch.
\end{itemize}

\subsection*{Bài học}
\begin{itemize}[nosep]
\item Phép đo trước/sau (snapshot) bắt được thứ test đơn vị không bắt được: hook quên \texttt{extra} cho ra trạng thái
  cuối giống hệt, chỉ khác id menu (FR-A mutation M7).
\item Đo UAT chỉ-đọc \emph{trước} mỗi lần tách (giá trị noupdate, số kế sequence, người trong nhóm) rẻ hơn nhiều so với
  sửa sau deploy.
\item \texttt{createdb -T} không chép filestore; bundle \texttt{/web/assets} trả 500 và làm sai mọi số đo browser.
\item Test quét component là hàng rào thật: CSS riêng một màn bị chặn, cách đúng là đổi bố cục bằng hàng chuẩn.
\end{itemize}

\subsection*{Nợ còn lại (chỉ ghi)}
\begin{itemize}[nosep]
\item \texttt{wujia\_portal\_sale} còn write ở route update/remove/note; \texttt{portal\_return} còn \texttt{write} gắn đính kèm.
\item \texttt{\_sql\_constraints} (Odoo 19 bỏ) ở 16 file, 8 phía Dev: chuyển \texttt{models.Constraint} khi có phiên deprecations.
\item Comment dẫn mã issue: 49 dòng \texttt{.py} / 349 dòng mọi file ngoài khối A, dọn dần khi chạm màn.
\item \texttt{auto\_install} theo ADR-027 đã chốt (chapter 74) nhưng chưa áp code; áp đồng loạt một phiên.
\item Bàn giao anh Thái: R7 $\times$2 \texttt{\_wj\_ensure\_contract} · 3 module mobile \texttt{auto\_install} · tests
  \texttt{wujia\_franchise} import file đã xoá · 8 file \texttt{\_sql\_constraints}.
\item Câu hỏi BA: 7 câu ADR-027 + câu tồn từ F9--F13 (phạm vi ticket theo người tạo, KPI đổi trả có tính ``đang xét'',
  nhóm ``User'' của Thi/Đổi trả/Thông báo đang 0 người).
\end{itemize}

\textbf{Sau khối A}: Issue List lại là việc chính. 8 issue mở chia 4 cụm: 140+141 shell header/cửa hàng · 143+144+145
mật độ mobile · 142 Home PC · 146 routing. Theo review FR-A, trong các component chưa có issue chỉ EmptyState đáng mở
issue mới.
